\chapter{Project Design}
\label{ch:design}

\section{Project Design Overview}
This chapter presents the design of the proposed Bangla propaganda detection
system. It describes the system requirements, overall interaction and data flow,
user-interface design, the detailed methodology (including alternative solutions
considered and the reasons for selecting the proposed approach), the project plan,
and the allocation of tasks among team members.

\section{Requirement Analysis}
Requirement analysis identifies what the proposed system should do and what
constraints it should satisfy. The requirements are divided into functional and
nonfunctional requirements, followed by diagrams showing how users, data, and
system components interact.

\subsection{Functional and Nonfunctional Requirements}

\subsubsection*{Functional Requirements}
\begin{itemize}
  \item \textbf{FR1: User Input.} The system shall allow a user to enter or submit
        Bangla text for analysis.
  \item \textbf{FR2: Text Preprocessing.} The system shall clean and prepare the
        submitted Bangla text before sending it to the detection model.
  \item \textbf{FR3: Propaganda Detection.} The system shall analyze the input text
        and determine whether propaganda is present.
  \item \textbf{FR4: Propaganda Technique Identification.} When propaganda is
        detected, the system should identify the possible propaganda technique,
        such as loaded language, fear appeal, or bandwagon, based on the final
        label set used in the dataset.
  \item \textbf{FR5: Result Display.} The system shall display the prediction
        clearly to the user (Input: Bangla social-media text; Prediction:
        Propaganda Detected; Technique: Loaded Language).
  \item \textbf{FR6: Confidence Information.} The system may display a confidence
        score or confidence level associated with the prediction.
  \item \textbf{FR7: Dataset Management.} The research component of the system
        shall support the collection, cleaning, annotation, verification, and
        storage of Bangla propaganda examples.
  \item \textbf{FR8: LLM-Assisted Pre-Annotation.} The dataset-development
        pipeline shall allow an LLM to generate an initial label suggestion for
        collected Bangla content.
  \item \textbf{FR9: Human Verification.} The system or research workflow shall
        allow a human annotator to review and correct the LLM's suggested label.
  \item \textbf{FR10: Model Training.} The research pipeline shall support training
        and fine-tuning of the selected propaganda detection model using the
        human-verified dataset.
  \item \textbf{FR11: Model Evaluation.} The system shall calculate evaluation
        measures including accuracy, precision, recall, and F1-score.
  \item \textbf{FR12: Iterative Dataset Improvement.} The research workflow shall
        identify poorly performing classes and allow additional data to be
        collected or generated and verified.
\end{itemize}

\subsubsection*{Nonfunctional Requirements}
\begin{itemize}
  \item \textbf{NFR1: Usability.} The user interface should be simple enough that a
        user can submit Bangla text and understand the prediction without
        technical knowledge.
  \item \textbf{NFR2: Performance.} The system should provide predictions within a
        reasonable response time for normal user requests.
  \item \textbf{NFR3: Reliability.} The system should provide consistent results
        when the same input is analyzed under the same model configuration.
  \item \textbf{NFR4: Accuracy.} The model should achieve acceptable performance
        according to the selected evaluation metrics.
  \item \textbf{NFR5: Maintainability.} The system should be designed so that the
        detection model and dataset can be updated without rebuilding the complete
        application.
  \item \textbf{NFR6: Scalability.} The architecture should allow the dataset and
        model to grow beyond the initial 500--1,000 custom samples.
  \item \textbf{NFR7: Security.} User-submitted data and research data should be
        handled securely and unnecessary personal information should not be stored.
  \item \textbf{NFR8: Explainability.} Where feasible, the system should provide
        useful information about why a text was classified as potentially
        propagandistic rather than displaying only a class label.
  \item \textbf{NFR9: Language Support.} The main detection workflow shall support
        Bangla text.
\end{itemize}

% ---------------------------------------------------------------------
\subsection{Context Diagram}
The context diagram represents the proposed system as a single high-level process
and shows its relationship with external entities. The main external entities are:
\begin{enumerate}
  \item \textbf{User:} \emph{User $\rightarrow$ System:} Bangla text.
        \emph{System $\rightarrow$ User:} detection result, propaganda category,
        and optional confidence information.
  \item \textbf{Researcher / Administrator:} manages the research dataset and
        model-development process. \emph{Researcher $\rightarrow$ System:} raw
        data, annotation decisions, training configuration.
        \emph{System $\rightarrow$ Researcher:} dataset information, model
        results, evaluation metrics.
  \item \textbf{LLM Service:} used during dataset development as a pre-annotation
        assistant. \emph{System $\rightarrow$ LLM:} Bangla sample and annotation
        prompt. \emph{LLM $\rightarrow$ System:} suggested label, possible
        technique, and explanation. (The LLM output is \emph{not} considered the
        final ground truth; human verification is strictly required.)
  \item \textbf{External Data Sources:} publicly available datasets and other
        permitted sources provide the initial and additional Bangla data used
        during development. \emph{External Data Sources $\rightarrow$ System:}
        Bangla text/images and associated metadata where available.
\end{enumerate}

\begin{figure}[H]
\centering
\begin{adjustbox}{max width=\textwidth}
\begin{tikzpicture}
% central system
\node[draw=boxdraw,thick,fill=boxfill,rounded corners=6pt,minimum width=32mm,
      minimum height=66mm,text width=26mm,align=center,font=\footnotesize\bfseries] (S) at (0,0)
      {Bangla Propaganda Detection System};
% external entities
\node[ext] (U) at (-5.6, 2.2) {User};
\node[ext] (R) at (-5.6,-2.2) {Researcher / Administrator};
\node[ext] (L) at ( 5.6, 2.2) {LLM Service};
\node[ext] (E) at ( 5.6,-2.2) {External Data Sources};

% User <-> System
\draw[flow] (-4.29, 2.65) -- node[lab,above,text width=25mm]{Bangla text} (-1.6, 2.65);
\draw[flow] (-1.6, 1.75) -- node[lab,below,text width=25mm]{Detection result, category \& confidence} (-4.29, 1.75);
% Researcher <-> System
\draw[flow] (-4.29,-1.75) -- node[lab,above,text width=25mm]{Raw data, annotation decisions, training config} (-1.6,-1.75);
\draw[flow] (-1.6,-2.65) -- node[lab,below,text width=25mm]{Dataset info, model results, metrics} (-4.29,-2.65);
% System <-> LLM
\draw[flow] ( 1.6, 2.65) -- node[lab,above,text width=25mm]{Cleaned sample \& annotation prompt} ( 4.29, 2.65);
\draw[flow] ( 4.29, 1.75) -- node[lab,below,text width=25mm]{Suggested label, technique \& explanation} ( 1.6, 1.75);
% External data -> System
\draw[flow] ( 4.29,-2.2) -- node[lab,above,text width=25mm]{Raw Bangla text/images \& metadata} ( 1.6,-2.2);
\end{tikzpicture}
\end{adjustbox}
\caption{Context diagram --- the proposed system acts as the central component
connecting the user, research workflow, external datasets, and the LLM-based
annotation assistant.}
\label{fig:context}
\end{figure}

% ---------------------------------------------------------------------
\subsection{Data Flow Diagram}
\label{sec:dfd}

\subsubsection{Level 1 Data Flow Diagram}
The Level~1 Data Flow Diagram (DFD) expands the main system into its major internal
processes and data stores.

\paragraph{Processes.}
\begin{itemize}
  \item \textbf{P1: Data Collection and Management.} Collects Bangla data from
        existing datasets and permitted online sources.
  \item \textbf{P2: Data Preprocessing.} Cleans, normalizes, and prepares the
        collected data.
  \item \textbf{P3: LLM Pre-Annotation.} Sends the cleaned samples to the LLM and
        receives possible propaganda labels and explanations.
  \item \textbf{P4: Human Verification.} Allows the researcher or annotator to
        review the LLM output and produce the final verified label.
  \item \textbf{P5: Model Training.} Uses the verified dataset to fine-tune the
        selected machine-learning model.
  \item \textbf{P6: Model Evaluation.} Evaluates the trained model using accuracy,
        precision, recall, F1-score, and class-wise performance.
  \item \textbf{P7: User Prediction.} Receives new Bangla text from the user,
        preprocesses it, sends it to the trained model, and returns the prediction.
\end{itemize}

\paragraph{Main data stores.}
\begin{itemize}
  \item \textbf{D1: Raw Dataset} --- collected Bangla data before processing.
  \item \textbf{D2: Processed Dataset} --- cleaned and prepared data.
  \item \textbf{D3: Human-Verified Dataset} --- final dataset with human-confirmed
        labels.
  \item \textbf{D4: Trained Model} --- the trained propaganda detection model.
  \item \textbf{D5: Evaluation Results} --- performance measures and experimental
        results.
\end{itemize}

\begin{figure}[H]
\centering
\begin{adjustbox}{max width=\textwidth,max height=0.84\textheight}
\begin{tikzpicture}
% external entities (left)
\node[ext]  (eds) at (-5.3,  0.0) {External Data Sources};
\node[ext]  (llm) at (-5.3, -4.8) {LLM Service};
\node[ext]  (res) at (-5.3, -7.2) {Researcher / Expert};
\node[ext]  (usr) at (-5.3,-15.9) {End User};
% processes (centre)
\node[proc,text width=30mm] (p1) at (0,  0.0) {P1: Data Collection \& Management};
\node[proc,text width=30mm] (p2) at (0, -2.4) {P2: Data Preprocessing};
\node[proc,text width=30mm] (p3) at (0, -4.8) {P3: LLM Pre-Annotation};
\node[proc,text width=30mm] (p4) at (0, -7.2) {P4: Human Verification};
\node[proc,text width=30mm] (p5) at (0, -9.6) {P5: Model Training};
\node[proc,text width=30mm] (p6) at (0,-12.9) {P6: Model Evaluation (uses test split of D3)};
\node[proc,text width=30mm] (p7) at (0,-15.9) {P7: User Prediction};
% data stores (right)
\node[store,text width=26mm] (d1) at (5.3, -1.2) {D1: Raw Dataset};
\node[store,text width=26mm] (d2) at (5.3, -3.6) {D2: Processed Dataset};
\node[store,text width=26mm] (d3) at (5.3, -8.4) {D3: Human-Verified Dataset};
\node[store,text width=26mm] (d4) at (5.3,-11.0) {D4: Trained Model (weights \& config)};
\node[store,text width=26mm] (d5) at (5.3,-12.9) {D5: Evaluation Results};
\node[ext]   (r2) at (5.3,-14.9) {Researcher / Expert};

% ---- flows
\draw[flow] (eds) -- node[lab,above]{Raw content} (p1);
\draw[flow] (p1.east) -| (d1.north);
\draw[flow] (d1.west) -| (p2.north);
\draw[flow] (p2.east) -| (d2.north);
\draw[flow] (d2.west) -| (p3.north);
\draw[flow,<->] (p3.west) -- node[lab,above,text width=20mm]{Prompts \& suggestions} (llm.east);
\draw[flow] (p3.south) -- node[lab,right]{Pre-annotated} (p4.north);
\draw[flow,<->] (p4.west) -- node[lab,above,text width=20mm]{Review \& correct} (res.east);
\draw[flow] (p4.east) -| (d3.north);
\draw[flow] (d3.west) -| (p5.north);
\draw[flow] (p5.east) -| (d4.north);
\draw[flow] (d4.west) -| (p6.north);
\draw[flow] (p6.east) -- (d5.west);
\draw[flow] (d5.south) -- node[lab,right]{Results} (r2.north);
\draw[flow] (d4.east) -- ++(0.9,0) |- (p7.east);
\node[lab,above,text width=16mm] at (2.8,-15.9) {Inference weights};
% user <-> P7
\draw[flow] ($(usr.east)+(0,0.25)$) -- node[lab,above]{Bangla text} ($(p7.west)+(0,0.25)$);
\draw[flow] ($(p7.west)-(0,0.25)$) -- node[lab,below,text width=22mm]{Prediction, category \& explanation} ($(usr.east)-(0,0.25)$);
\end{tikzpicture}
\end{adjustbox}
\caption{Level~1 Data Flow Diagram (DFD) --- the dual-workflow architecture: the
research and model-development pipeline and the end-user prediction workflow.}
\label{fig:dfd1}
\end{figure}

\subsubsection{Detailed Level 2 Data Flow Diagram}
\label{sec:dfd2}
Figure~\ref{fig:dfd2} expands the Level~1 processes to show their internal
processing. Processes~1 and~2 harvest and normalize both text and images.
Processes~3 and~4 form the hybrid annotation engine, where few-shot prompts drive an
LLM to produce candidate suggestions that are placed in a buffer and then verified
by human annotators in a dedicated workspace before being written to the verified
ground-truth store (D3). Process~5 tokenizes verified text, encodes it with a
BanglaBERT/XLM-R backbone, encodes images with a ViT/ResNet extractor, and joins the
two through a cross-modal attention fusion layer feeding the propaganda technique
classification head. Processes~6 and~7 serve live user input through the inference
engine and pair each prediction with a feature-attribution explainability module.

\begin{figure}[H]
\centering
\begin{adjustbox}{max width=\textwidth,max height=0.86\textheight}
\begin{tikzpicture}
% ===== Processes 1 & 2
\node[blk]  (raw)   at (-5.25, 0.0) {Raw Online Bangla Content};
\node[proc] (scr)   at (-1.75, 0.0) {Automated Scraper \& API Harvester};
\node[proc] (img)   at ( 1.75, 0.0) {Image Preprocessing: Normalization \& Resizing};
\node[proc] (clean) at ( 1.75,-1.9) {Bangla Text Cleaner: Emoji, Punctuation, Normalizer};
\draw[grp] (-6.8,0.95) rectangle (6.8,-2.85);
\node[grptitle,anchor=south west] at (-6.8,0.95) {Processes 1 \& 2: Ingestion \& Preprocessing};
\draw[flow] (raw) -- (scr);
\draw[flow] (scr) -- (img);
\draw[flow] (scr.south) |- (clean.west);

% ===== Processes 3 & 4
\node[proc]  (prompt) at ( 1.75,-4.5) {Few-Shot Prompt Constructor};
\node[proc]  (llm)    at (-1.75,-4.5) {LLM Pre-Annotation API};
\node[store] (buf)    at (-5.25,-4.5) {Candidate Suggestion Buffer};
\node[proc]  (ui)     at (-5.25,-6.4) {Human Annotator Verification Workspace};
\node[store] (gt)     at (-1.75,-6.4) {D3: Verified Ground-Truth Store};
\draw[grp] (-6.8,-3.55) rectangle (6.8,-7.35);
\node[grptitle,anchor=south west] at (-6.8,-3.55) {Processes 3 \& 4: Hybrid Annotation Engine};
\draw[flow] (clean.south) -- (prompt.north);
\draw[flow] (prompt) -- (llm);
\draw[flow] (llm) -- (buf);
\draw[flow] (buf) -- (ui);
\draw[flow] (ui) -- (gt);

% ===== Process 5
\node[proc] (tok)  at (-1.75, -9.0) {Bangla WordPiece / SentencePiece Tokenizer};
\node[proc] (bb)   at ( 1.75, -9.0) {BanglaBERT / XLM-R Transformer Backbone};
\node[proc] (fuse) at ( 1.75,-10.9) {Cross-Modal Attention Fusion Layer};
\node[proc] (vit)  at ( 5.25,-10.9) {Vision Transformer / ResNet Visual Extractor};
\node[proc] (cls)  at (-1.75,-10.9) {Propaganda Technique Classification Head};
\draw[grp] (-6.8,-8.05) rectangle (6.8,-11.85);
\node[grptitle,anchor=south east] at (6.8,-8.05) {Process 5: Feature Extraction \& Model Training};
\draw[flow] (gt.south) -- (tok.north);
\draw[flow] (tok) -- (bb);
\draw[flow] (bb) -- (fuse);
\draw[flow] (vit) -- (fuse);
\draw[flow] (fuse) -- (cls);
% image path (outside right)
\draw[dflow] (img.east) -- (7.05,0) |- (vit.east);
\node[lab,rotate=90,anchor=north] at (7.1,-5.5) {Image features};

% ===== Processes 6 & 7
\node[ext]  (uin)  at (-5.25,-14.4) {User Bangla Text / URL};
\node[proc] (inf)  at (-1.75,-14.4) {Model Inference Engine};
\node[proc] (pout) at ( 1.75,-13.5) {Propaganda Label \& Technique Identification};
\node[proc] (shap) at ( 1.75,-15.3) {Feature Attribution \& Explainability Module};
\node[proc] (disp) at ( 5.25,-14.4) {Web Results Dashboard};
\draw[grp] (-6.8,-12.55) rectangle (6.8,-16.25);
\node[grptitle,anchor=south east] at (6.8,-12.55) {Processes 6 \& 7: Prediction \& Explainability};
\draw[flow] (uin) -- (inf);
\draw[flow] (inf.east)  -- ++(0.43,0) |- (pout.west);
\draw[flow] (inf.east)  -- ++(0.43,0) |- (shap.west);
\draw[flow] (pout.east) -- ++(0.43,0) |- (disp.west);
\draw[flow] (shap.east) -- ++(0.43,0) |- (disp.west);
\draw[dflow] (cls.south) -- (inf.north);
\node[lab,anchor=west] at (-1.6,-12.2) {Trained weights};
\end{tikzpicture}
\end{adjustbox}
\caption{Detailed Level~2 Data Flow Diagram (DFD) for the proposed system.}
\label{fig:dfd2}
\end{figure}

% ---------------------------------------------------------------------
\subsection{UI Design}
The proposed application uses a simple, intuitive interface so that users can
easily submit Bangla content and understand the result.

\subsubsection*{Main Screen}
The main screen contains the application title (\textbf{Bangla Propaganda Detection
System}), a short explanation of the system purpose, a Bangla text input area, an
``Analyze'' button, and a result section.

\begin{figure}[H]
\centering
\begin{adjustbox}{max width=\textwidth}
\begin{tikzpicture}[ui/.style={font=\footnotesize,anchor=west}]
\draw[boxdraw,thick,rounded corners=4pt,fill=white] (-6.2,0) rectangle (6.2,-9.4);
\node[font=\large\bfseries] at (0,-0.75) {Bangla Propaganda Detection System};
\node[font=\footnotesize] at (0,-1.4) {Analyze Bangla digital content for persuasive \& propaganda cues};
\draw[boxdraw] (-5.8,-1.85) -- (5.8,-1.85);

\node[ui,font=\footnotesize\bfseries] at (-5.8,-2.4) {Enter Bangla Text:};
\draw[gray,rounded corners=2pt,fill=gray!6] (-5.8,-2.75) rectangle (5.8,-4.55);
\node[ui,text=gray] at (-5.6,-3.15) {\bn{এখানে বাংলা সোশ্যাল মিডিয়া পোস্ট বা খবর লিখুন\ldots}};

\node[draw=boxdraw,fill=uiublue,text=white,rounded corners=3pt,minimum width=32mm,
      minimum height=8mm,font=\footnotesize\bfseries] at (0,-5.3) {Analyze};
\draw[boxdraw] (-5.8,-6.0) -- (5.8,-6.0);

\node[ui,font=\footnotesize\bfseries] at (-5.8,-6.5) {Result:};
\node[ui,anchor=north west,text width=11cm,align=left] at (-5.6,-6.9)
  {\textbf{Prediction:} Propaganda Detected\\[2pt]
   \textbf{Technique:} Loaded Language\\[2pt]
   \textbf{Confidence:} Medium\\[2pt]
   \textbf{Explanation:} The text contains emotionally strong wording that may influence the reader.};
\end{tikzpicture}
\end{adjustbox}
\caption{Main screen result interface.}
\label{fig:ui-main}
\end{figure}

\subsubsection*{Result Interface}
The result page clearly separates the input, prediction, and explanation:
\textbf{Prediction:} Propaganda Detected; \textbf{Possible Technique:} Loaded
Language; \textbf{Confidence:} Medium; \textbf{Explanation:} the text contains
emotionally strong wording that may influence the reader. The explanation is
treated as supporting information rather than as a replacement for the
classification result.

\subsubsection*{Annotation Interface}
A separate interface is used by the researcher or annotator to support the LLM
pre-annotation $\rightarrow$ human verification workflow.

\begin{figure}[H]
\centering
\begin{adjustbox}{max width=\textwidth}
\begin{tikzpicture}[ui/.style={font=\footnotesize,anchor=west},
                    b/.style={ui,font=\footnotesize\bfseries}]
\draw[boxdraw,thick,rounded corners=4pt,fill=white] (-6.2,0) rectangle (6.2,-11.0);
\node[font=\large\bfseries] at (0,-0.7) {Researcher / Annotator Interface};
\draw[boxdraw] (-5.8,-1.2) -- (5.8,-1.2);

\node[ui] at (-5.8,-1.7) {\textbf{Sample ID:} \#1042};
\node[ui] at (-5.8,-2.3) {\textbf{Text:} ``\ldots'' (Bangla sample text)};

\node[b] at (-5.8,-3.1) {LLM Suggested Label:};
\node[ui,draw=boxdraw,fill=boxfill,inner sep=3pt] at (-2.2,-3.1) {Propaganda Detected};
\node[b] at (-5.8,-3.8) {Suggested Technique:};
\node[ui,draw=boxdraw,fill=boxfill,inner sep=3pt] at (-2.2,-3.8) {Fear Appeal};
\node[b] at (-5.8,-4.5) {LLM Rationale:};
\node[ui] at (-2.2,-4.5) {Mentions severe imminent societal threat};

\node[b] at (-5.8,-5.4) {Human Decision:};
\draw (-5.5,-6.0) circle (0.11);  \node[ui] at (-5.25,-6.0) {Accept LLM Suggestion};
\draw (-5.5,-6.65) circle (0.11); \node[ui] (mod) at (-5.25,-6.65) {Modify Label:};
\node[ui,draw=gray,inner sep=3pt,right=3mm of mod] {Select Technique $\blacktriangledown$};
\draw (-5.5,-7.3) circle (0.11);  \node[ui] at (-5.25,-7.3) {Mark Ambiguous / Sarcastic};

\node[b] at (-5.8,-8.05) {Annotator Notes:};
\draw[gray,rounded corners=2pt,fill=gray!6] (-5.8,-8.35) rectangle (5.8,-9.35);

\node[draw=boxdraw,fill=uiublue,text=white,rounded corners=3pt,minimum width=48mm,
      minimum height=8mm,font=\footnotesize\bfseries] at (0,-10.2) {Confirm Ground Truth};
\end{tikzpicture}
\end{adjustbox}
\caption{Annotation interface.}
\label{fig:ui-annot}
\end{figure}

% ---------------------------------------------------------------------
\section{Detailed Methodology and Design}
\label{sec:detailed-method}
This section explains how the proposed system will be developed and why the
selected solutions are appropriate. Several alternative approaches were considered
before selecting the proposed combination of existing datasets, human-verified
custom data, transformer-based models, and LLM-assisted annotation.

\subsection{Overall Design Approach}
The project follows an iterative design rather than attempting to build the
complete system in a single step:
\[
\begin{gathered}
\text{Existing Dataset} \rightarrow \text{Baseline Pipeline} \rightarrow \text{Custom Dataset} \\
\rightarrow \text{Annotation} \rightarrow \text{Model Training} \rightarrow \text{Evaluation} \rightarrow \text{Improvement}
\end{gathered}
\]
This approach reduces development risk because the model-development pipeline can
be tested before the custom dataset is complete.

\subsection{Alternative Solutions Considered}

\subsubsection{Alternative 1: Building a Large Dataset Completely From Scratch}
One possible approach was to immediately collect and manually annotate a very large
Bangla dataset (e.g., attempting to label 10,000--50,000 samples before beginning
model development).
\begin{itemize}
  \item \textbf{Advantages:} Labels would be created specifically for the research
        task; the dataset could be designed around the exact propaganda categories
        required.
  \item \textbf{Disadvantages:} Manual annotation would require a large amount of
        time; annotation cost would increase significantly; difficult and
        ambiguous examples would require additional review; model development
        would need to wait until a large dataset was completed.
  \item \textbf{Selected Approach:} The proposed research instead begins with
        existing datasets and gradually creates a smaller custom dataset. This
        allows the technical pipeline to be developed immediately and reduces the
        initial annotation burden.
\end{itemize}

\subsubsection{Alternative 2: Fully Automatic LLM Annotation}
Another possible solution was to ask an LLM to label all collected samples and
directly use those labels as the dataset ground truth.
\begin{itemize}
  \item \textbf{Advantages:} Very fast compared with manual annotation; large
        quantities of data can be labeled automatically; reduces human workload.
  \item \textbf{Disadvantages:} LLM predictions can be incorrect; ambiguous Bangla
        language may be interpreted incorrectly; the LLM may introduce systematic
        labeling bias; using LLM labels as ground truth may reduce dataset
        reliability.
  \item \textbf{Selected Approach:} The proposed system uses the LLM only for
        pre-annotation. Every final label must be checked by a human annotator.
        Therefore
        \[\text{LLM Suggestion}\neq\text{Ground Truth},\]
        and the ground truth is the final human-verified label.
\end{itemize}

\subsubsection{Alternative 3: Complete Manual Annotation}
A third approach is to manually annotate every collected sample without using an LLM.
\begin{itemize}
  \item \textbf{Advantages:} Direct human judgment; better control over the
        annotation process; no dependence on LLM-generated suggestions.
  \item \textbf{Disadvantages:} Very time-consuming; more expensive in terms of
        researcher effort; difficult to scale as the dataset grows; annotators may
        become tired, which can reduce consistency.
  \item \textbf{Selected Approach:} The proposed LLM-assisted approach provides a
        compromise. The LLM performs an initial classification, while humans make
        the final decision. This reduces repetitive work without removing human
        control.
\end{itemize}

\subsubsection{Alternative 4: Traditional Machine Learning}
Traditional models such as SVM, Logistic Regression, or Naive Bayes could be used as
the main final classifier.
\begin{itemize}
  \item \textbf{Advantages:} Simple to implement; fast to train; computationally
        inexpensive; useful as baseline models.
  \item \textbf{Disadvantages:} Usually require manually designed text features;
        may have difficulty capturing complex contextual relationships in Bangla;
        may perform poorly when propaganda depends on wider context.
  \item \textbf{Selected Approach:} A traditional model is still implemented as a
        baseline. However, the main experiments investigate transformer-based
        models because they capture contextual language information more
        effectively.
\end{itemize}

\subsubsection{Alternative 5: BanglaBERT vs.\ XLM-R}
Two candidate transformer approaches are particularly relevant:
\begin{itemize}
  \item \textbf{BanglaBERT.} \emph{Advantages:} designed specifically for Bangla;
        suitable for Bangla text classification; can provide strong
        language-specific representations. \emph{Disadvantages:} primarily focused
        on Bangla; may be less suitable if the project is later expanded to
        multiple languages.
  \item \textbf{XLM-R.} \emph{Advantages:} supports many languages; can support
        multilingual expansion; useful if the system later needs to process both
        Bangla and other languages. \emph{Disadvantages:} not designed only for
        Bangla; a Bangla-specific model may be more suitable for some Bangla-only
        tasks.
  \item \textbf{Selected Approach:} BanglaBERT is used as the primary candidate and
        XLM-R as a comparison model. This allows the research to determine whether
        a Bangla-specific model or a multilingual model performs better on the
        custom dataset instead of assuming the answer beforehand.
\end{itemize}

\subsubsection{Alternative 6: Text-Only vs.\ Multimodal Detection}
A text-only system is simpler because it requires only Bangla text. However,
propaganda can also appear through images, memes, and combinations of text and
visual content. Existing research shows that multimodal information can contain
useful detection signals.
\begin{itemize}
  \item \textbf{Text-Only Approach.} \emph{Advantages:} easier to implement;
        requires fewer computational resources; easier dataset collection;
        suitable for the initial prototype. \emph{Disadvantages:} cannot analyze
        visual propaganda; may miss important information in images.
  \item \textbf{Multimodal Approach.} \emph{Advantages:} uses both text and image
        information; more suitable for posts containing memes or image-based
        messages. \emph{Disadvantages:} more complex; requires image data and
        suitable multimodal models; requires more computational resources;
        annotation becomes more difficult.
  \item \textbf{Selected Approach:} The initial system focuses on Bangla text
        detection so that the core pipeline can be developed and evaluated
        properly. Multimodal support is investigated as an extension where
        suitable image--text data are available. This follows the project's
        iterative strategy: first make the core system reliable, then increase its
        complexity.
\end{itemize}

\subsection{Design Rationale and Trade-off Evaluation}
Table~\ref{tab:tradeoff} consolidates the alternatives above into a single
evaluation matrix, stating for each design layer what was considered and why the
selected option was chosen.

\begingroup
\small
\renewcommand{\arraystretch}{1.3}
\begin{longtable}{@{}L{2.9cm}L{3.9cm}L{6.1cm}@{}}
\caption{Design rationale and trade-off evaluation matrix.}\label{tab:tradeoff}\\
\toprule
\textbf{Design Layer} & \textbf{Considered Alternative} & \textbf{Selected Choice \& Engineering Rationale}\\
\midrule
\endfirsthead
\multicolumn{3}{@{}l}{\small\itshape Table~\thetable\ (continued)}\\
\toprule
\textbf{Design Layer} & \textbf{Considered Alternative} & \textbf{Selected Choice \& Engineering Rationale}\\
\midrule
\endhead
\bottomrule
\endlastfoot

\textbf{Dataset Creation Strategy} & Fully manual annotation from scratch (10,000+ samples). &
\textbf{Hybrid LLM pre-annotation + human verification:} reduces manual annotator fatigue and labeling cost while preserving ground-truth fidelity through expert human sign-off~\cite{maarouf2024}.\\
\addlinespace
\textbf{Annotation Authority} & Fully autonomous LLM zero-shot labeling. &
\textbf{Human-in-the-loop ground truth:} pure LLM labels can suffer from hallucination and cultural blindness in Bangla idioms; human verification protects label quality~\cite{hasan2026,maarouf2024}.\\
\addlinespace
\textbf{Core NLP Backbone} & Classical machine learning (TF-IDF + SVM / Naive Bayes). &
\textbf{Fine-tuned pretrained transformers (BanglaBERT \& XLM-R):} capture bidirectional contextual semantics, sarcasm, and rhetorical nuance that $n$-gram models miss. Classical ML is retained strictly as the E1 baseline.\\
\addlinespace
\textbf{Monolingual vs.\ Multilingual} & Generic multilingual models only (mBERT / XLM-R). &
\textbf{Dual evaluation (BanglaBERT vs.\ XLM-R):} BanglaBERT provides vocabulary specialization for Bengali script, while XLM-R tests cross-lingual robustness against code-mixed Bangla--English text.\\
\addlinespace
\textbf{Modality Scope} & Full multimodal (video + audio + text + network) from day one. &
\textbf{Phased, extensible modular architecture:} Phase~1 establishes robust text detection and dataset verification; Phase~2 adds image--text meme fusion; Phase~3 integrates graph-based coordination metrics.\\
\addlinespace
\textbf{Class Imbalance Strategy} & Standard cross-entropy on the raw class distribution. &
\textbf{Cost-sensitive loss weighting + strategic oversampling:} prevents bias toward the dominant ``No Propaganda'' class and supports viable recall on minority techniques (e.g., Bandwagon, Fear Appeal).\\
\end{longtable}
\endgroup

\subsection{Dataset Design}
The dataset design follows a three-stage methodology:
\begin{itemize}
  \item \textbf{Stage 1: Existing Dataset.} Existing publicly available Bangla
        datasets are used to develop and test the initial pipeline.
  \item \textbf{Stage 2: Custom Dataset.} Approximately 500--1,000 Bangla samples
        will initially be collected from appropriate public sources.
  \item \textbf{Stage 3: Dataset Improvement.} After model evaluation,
        underrepresented classes will be identified and additional real examples
        will be collected where necessary. Synthetic augmentation may be
        investigated only when real examples remain insufficient, following the
        workflow
        \[\text{Generate}\rightarrow\text{Filter}\rightarrow\text{Human Check}\rightarrow\text{Training Use}.\]
        Synthetic examples will be stored separately from real-world examples
        during evaluation.
\end{itemize}

\subsection{Annotation Design}
The annotation system is based on a two-level process:
\begin{itemize}
  \item \textbf{Level 1: LLM Pre-Annotation.} The LLM receives a Bangla sample and
        suggests whether propaganda may be present, the possible propaganda
        technique, confidence or certainty, and a short explanation.
  \item \textbf{Level 2: Human Verification.} The human annotator checks the
        prediction and assigns the final label.
\end{itemize}
The final dataset therefore contains
\[\text{Original Text}+\text{LLM Suggestion}+\text{Human Decision}+\text{Final Label}.\]
This information can also help researchers analyze where the LLM performs well or
poorly.

\subsection{Class Imbalance Design}
The dataset may contain significantly different numbers of examples for different
propaganda categories. For example: \emph{No Propaganda}: 700; \emph{Loaded
Language}: 100; \emph{Fear Appeal}: 40; \emph{Bandwagon}: 20; \emph{Other}: 10.
The project will monitor class distributions after each data-collection cycle.
Possible solutions include:
\begin{enumerate}
  \item \textbf{Additional real-data collection:} more minority-class examples will
        be collected where possible.
  \item \textbf{Class weighting:} the training algorithm may assign greater
        importance to minority classes.
  \item \textbf{Oversampling:} minority examples may be sampled more frequently
        during training.
  \item \textbf{Synthetic augmentation:} LLM-generated examples may be investigated
        after human filtering.
\end{enumerate}
The chosen technique will be determined experimentally rather than assumed to be
effective in advance.

\subsection{Training and Testing Design}
\begin{itemize}
  \item The human-verified dataset will be divided into \textbf{training},
        \textbf{validation}, and \textbf{test} subsets.
  \item The training set will be used for model learning.
  \item The validation set will be used for model selection and parameter tuning.
  \item The test set will be kept separate and used for final evaluation.
  \item The test data must not be used during training or annotation decisions in a
        way that causes data leakage.
\end{itemize}

\subsection{Experimental Design}
The project will compare multiple approaches rather than reporting only one model.
The experimental structure is shown in Table~\ref{tab:experiments}.

\begin{table}[H]
\centering
\caption{Experimental design.}\label{tab:experiments}
\renewcommand{\arraystretch}{1.25}
\begin{tabular}{@{}cl@{}}
\toprule
\textbf{Experiment} & \textbf{Approach}\\
\midrule
E1 & Traditional ML baseline\\
E2 & BanglaBERT\\
E3 & XLM-R\\
E4 & BanglaBERT + class-imbalance strategy\\
E5 & XLM-R + class-imbalance strategy\\
E6 & Optional LLM-assisted / synthetic augmentation experiment\\
\bottomrule
\end{tabular}
\end{table}
The final selection will be based on measured performance rather than assuming that
the most complex model will perform best.

\subsection{Evaluation Design}
The main performance measures will be:
\begin{itemize}
  \item \textbf{Accuracy:} overall proportion of correctly classified examples.
  \item \textbf{Precision:} how many predicted positive examples are actually positive.
  \item \textbf{Recall:} how many actual positive examples are correctly detected.
  \item \textbf{F1-score:} the harmonic mean of precision and recall.
\end{itemize}
For an imbalanced multi-class dataset, \textbf{macro-F1} and \textbf{class-wise
F1-score} will receive particular attention because they provide a clearer view of
minority-class performance. A confusion matrix will also be used to identify which
propaganda categories are commonly confused with one another.

\subsection{Error Analysis}
After evaluation, incorrectly classified examples will be manually inspected. The
researchers will investigate whether errors are caused by ambiguous language,
sarcasm, insufficient training examples, class imbalance, cultural context,
code-mixing, similar propaganda categories, or limitations of the selected model.
The error analysis will guide the next iteration of dataset collection and model
training.

\subsection{Final System Architecture}
The final proposed architecture is divided into two connected parts: a research
pipeline responsible for building and improving the model, and a user application
that uses the trained model to provide predictions (Figure~\ref{fig:arch}).

\begin{figure}[H]
\centering
\begin{adjustbox}{max width=\textwidth}
\begin{tikzpicture}
% ===== Research pipeline
\node[blk]   (ds)  at (-5.1, 0.0) {Data Sources};
\node[blk]   (dc)  at (-1.7, 0.0) {Data Collection};
\node[blk]   (dp)  at ( 1.7, 0.0) {Pre-processing};
\node[blk]   (lpa) at ( 5.1, 0.0) {LLM Pre-Annotation};
\node[blk]   (hv)  at ( 5.1,-1.9) {Human Verification};
\node[store] (vd)  at ( 1.7,-1.9) {Verified Dataset};
\node[blk]   (mt)  at (-1.7,-1.9) {Model Training};
\node[blk]   (ev)  at (-1.7,-3.8) {Evaluation};
\node[blk]   (mi)  at (-5.1,-3.8) {Model Improvement};
\draw[grp] (-6.8,0.95) rectangle (6.8,-4.75);
\node[grptitle,anchor=south west] at (-6.8,0.95) {Research Pipeline};
\draw[flow] (ds) -- (dc);  \draw[flow] (dc) -- (dp);  \draw[flow] (dp) -- (lpa);
\draw[flow] (lpa) -- (hv); \draw[flow] (hv) -- (vd);  \draw[flow] (vd) -- (mt);
\draw[flow] (mt) -- (ev);  \draw[flow] (ev) -- (mi);
\draw[dflow] (mi.north) -- (-5.1,-0.95) -| (dc.south);
\node[lab,anchor=west] at (-5.0,-2.3) {iterate};

% ===== User application
\node[ext]  (u)   at ( 5.1,-7.5) {User};
\node[blk]  (bti) at ( 1.7,-7.5) {Bangla Text Input};
\node[blk]  (up)  at (-1.7,-7.5) {Pre-processing};
\node[proc] (tpm) at (-5.1,-7.5) {Trained Propaganda Model};
\node[blk]  (pr)  at (-5.1,-9.4) {Prediction};
\node[blk]  (tce) at (-1.7,-9.4) {Technique / Confidence / Explanation};
\node[blk]  (res) at ( 1.7,-9.4) {User Result Display};
\draw[grp] (-6.8,-6.55) rectangle (6.8,-10.35);
\node[grptitle,anchor=south east] at (6.8,-6.55) {User Application};
\draw[flow] (u) -- (bti);  \draw[flow] (bti) -- (up);  \draw[flow] (up) -- (tpm);
\draw[flow] (tpm) -- (pr); \draw[flow] (pr) -- (tce);  \draw[flow] (tce) -- (res);

% deploy link
\draw[dflow] (ev.south) -- (-1.7,-5.6) -| (tpm.north);
\node[lab,above] at (-3.4,-5.6) {deploy model};
\end{tikzpicture}
\end{adjustbox}
\caption{Final system architecture: research pipeline and user application.}
\label{fig:arch}
\end{figure}

\subsection{Why the Proposed Design Was Selected}
The proposed design was selected because it balances research reliability,
development effort, scalability, and practicality. Using existing datasets allows
the technical pipeline to be developed early. Creating a smaller custom dataset
makes it possible to focus annotation effort on data specifically relevant to
Bangla propaganda detection. Using an LLM for pre-annotation reduces repetitive
work, while human verification protects the quality of the ground-truth labels.

Similarly, using traditional machine learning as a baseline and transformer models
as the main candidates allows meaningful comparison between simple and advanced
approaches. Starting with text-based detection keeps the initial implementation
manageable, while the architecture leaves room for future multimodal expansion.
The overall design is therefore iterative:
\[\text{Build}\rightarrow\text{Test}\rightarrow\text{Analyze}\rightarrow\text{Improve}\rightarrow\text{Retrain}\rightarrow\text{Re-evaluate}.\]
This makes it possible to improve both the dataset and the model throughout the
project instead of expecting the first version to be perfect.

% ---------------------------------------------------------------------
\section{Project Plan}
\label{sec:plan}
The project lifecycle is planned across three consecutive academic trimesters
(25 weeks), structured into four distinct engineering phases. The schedule is shown
as a Gantt chart in Figure~\ref{fig:gantt}.

\begin{figure}[H]
\centering
\begin{adjustbox}{max width=\textwidth}
\begin{ganttchart}[
    hgrid, vgrid={*{4}{draw=none},dotted},
    x unit=0.42cm, y unit chart=0.6cm, y unit title=0.55cm,
    title height=1, title label font=\tiny\bfseries,
    bar height=0.55, bar label font=\scriptsize,
    group label font=\scriptsize\bfseries,
    bar/.append style={fill=uiublue!70,draw=uiublue},
    group/.append style={fill=uiublue!35,draw=uiublue},
    bar label node/.append style={text width=5.8cm,align=right,anchor=east,font=\scriptsize\RaggedLeft},
    group label node/.append style={text width=5.8cm,align=right,anchor=east,font=\scriptsize\bfseries\RaggedLeft}
  ]{1}{25}
  \gantttitle{Weeks 1--25}{25}\\
  \gantttitlelist{1,...,25}{1}\\
  \ganttgroup{Phase 1: Planning}{1}{8}\\
  \ganttbar{Literature Review \& Taxonomy}{1}{8}\\
  \ganttbar{Requirements \& Architecture Spec.}{3}{8}\\
  \ganttgroup{Phase 2: Data Pipeline}{5}{16}\\
  \ganttbar{Web Scraping \& Data Ingestion}{5}{14}\\
  \ganttbar{LLM Pre-Annotation Prompting}{7}{12}\\
  \ganttbar{Human Expert Annotation \& Verification}{9}{16}\\
  \ganttgroup{Phase 3: Modeling \& Evaluation}{15}{22}\\
  \ganttbar{Classical Baselines (SVM/NB)}{15}{18}\\
  \ganttbar{BanglaBERT \& XLM-R Fine-Tuning}{16}{21}\\
  \ganttbar{Class Imbalance Mitigation \& Testing}{18}{22}\\
  \ganttgroup{Phase 4: UI \& Deployment}{20}{25}\\
  \ganttbar{Explainability Module \& UI Dashboard}{20}{24}\\
  \ganttbar{Benchmarking \& Final Thesis Draft}{22}{25}
\end{ganttchart}
\end{adjustbox}
\caption{Gantt chart of the project execution lifecycle (Weeks 1--25).}
\label{fig:gantt}
\end{figure}

\begin{itemize}
  \item \textbf{Phase 1: Requirements Gathering \& Taxonomy Definition (Weeks
        1--8).} Systematic review of propaganda literature, defining the Bangla
        propaganda technique taxonomy, and architecting the DFD.
  \item \textbf{Phase 2: Data Curation \& Hybrid Annotation (Weeks 5--16).}
        Gathering raw posts, prompt engineering for LLM pre-annotation, and
        conducting human expert verification to construct the verified dataset.
  \item \textbf{Phase 3: Model Training \& Experimental Benchmarking (Weeks
        15--22).} Fine-tuning BanglaBERT and XLM-R, implementing cost-sensitive
        loss weighting, and conducting ablation experiments E1--E6.
  \item \textbf{Phase 4: Dashboard Integration \& Final Thesis (Weeks 20--25).}
        Integrating the prediction engine with the web user interface, generating
        explainability summaries, and writing the final thesis report.
\end{itemize}

\section{Task Allocation}
\label{sec:tasks}
Project responsibilities are divided among team members according to technical
sub-disciplines:
\begin{itemize}
  \item \textbf{Shekh Abdullah Al Mehedi (ID: 112231075):} data pipeline
        architecture, automated web scraping infrastructure, data
        cleaning/normalization, and database management.
  \item \textbf{Mahjabin Khan (ID: 112230177):} annotation pipeline design, LLM
        prompt engineering for pre-annotation, annotator guideline drafting, and
        human verification coordination.
  \item \textbf{Omor Faruck Ullas (ID: 112310384):} machine learning pipeline,
        transformer model fine-tuning (BanglaBERT \& XLM-R), class imbalance
        mitigation, and loss function engineering.
  \item \textbf{Md.\ Rayan Rahman (ID: 112331071):} system evaluation, benchmarking
        experiments (E1--E6), confusion matrix error analysis, and web UI/API
        development.
\end{itemize}